\documentclass[11pt,a4paper]{article}
\usepackage[T1]{fontenc}
\usepackage[utf8]{inputenc}
\usepackage{lmodern}
\usepackage{microtype}
\usepackage[a4paper,margin=2.25cm]{geometry}
\usepackage{longtable,booktabs,array,calc}
\usepackage{fancyvrb}
\usepackage{xcolor}
\usepackage{hyperref}
\usepackage{xurl}
\usepackage{fancyhdr}
\usepackage{enumitem}
\usepackage{parskip}
\usepackage{titlesec}
\definecolor{ddtadark}{RGB}{30,38,48}
\definecolor{ddtablue}{RGB}{38,82,126}
\hypersetup{colorlinks=true,linkcolor=ddtablue,urlcolor=ddtablue,citecolor=ddtablue}
\pagestyle{fancy}
\fancyhf{}
\lhead{Documentation-Driven Threat Analysis}
\rhead{Research work plan R16}
\cfoot{\thepage}
\setlength{\headheight}{14pt}
\setlist[itemize]{leftmargin=1.5em,itemsep=0.2em,topsep=0.25em}
\setlist[enumerate]{leftmargin=1.7em,itemsep=0.2em,topsep=0.25em}
\titleformat{\section}{\Large\bfseries\color{ddtadark}}{}{0em}{}
\titleformat{\subsection}{\large\bfseries\color{ddtadark}}{}{0em}{}
\titleformat{\subsubsection}{\normalsize\bfseries\color{ddtadark}}{}{0em}{}
\providecommand{\tightlist}{\setlength{\itemsep}{0pt}\setlength{\parskip}{0pt}}
\setlength{\emergencystretch}{4em}
\hbadness=10000
\hfuzz=20pt
\sloppy
\begin{document}
\section{DDTA research work plan after documentation-layer closure}\label{ddta-research-work-plan-after-documentation-layer-closure}

\textbf{WORK PLAN - REVISION 16}

\textbf{Status:} active plan after BA3 integrated closure.

\textbf{Supersedes:} Revision 15 only for forward execution state; R1-R15 remain historical research records.

\subsection{1. Fixed prior state}\label{fixed-prior-state}

\begin{itemize}
\tightlist
\item
  Chapters 2-4: CLOSED / FINAL for current thesis scope.
\item
  Documentation layer: CLOSED.
\item
  BA0-R systems-modeling prior-art research: CLOSED.
\item
  BA0 responsibility/non-goals: CLOSED.
\item
  BA1 minimal BAE identity ontology: CLOSED.
\item
  BA2 relation/action vocabulary: CLOSED BY BA2-T4.
\item
  BA3 provenance/derivation/identity/lifecycle/change-revalidation mechanics: CLOSED BY BA3-T4.
\item
  \texttt{BAReferent} and \texttt{BAProposition}: ACCEPTED first-class semantic identity families.
\item
  ThreatForge is a case study/tooling instrument, not DDTA semantic authority.
\end{itemize}

\subsection{2. BA3 closure result}\label{ba3-closure-result}

BA3 closes the smallest current-scope contract for:

\begin{itemize}
\tightlist
\item
  independent baseline-scoped source provenance on both BA1 identity families;
\item
  \texttt{GROUNDED\ \textbar{}\ DERIVED\ \textbar{}\ DIAGNOSTIC\_UNRESOLVED} origin semantics;
\item
  role-bound derivation basis and immutable inspectable rule revisions;
\item
  \texttt{PENDING\_REVIEW\ \textbar{}\ ACCEPTED\ \textbar{}\ REJECTED} review state;
\item
  \texttt{CURRENT\ \textbar{}\ STALE} freshness;
\item
  family-specific cross-baseline semantic identity;
\item
  \texttt{RETAIN\ \textbar{}\ REPLACE\ \textbar{}\ RETIRE} continuity with historical \texttt{SUPERSEDED}/\texttt{RETIRED} interpretations;
\item
  explicit narrow \texttt{revalidationContext} and localized staleness propagation; and
\item
  the governed-document authority gate through corrective feedback.
\end{itemize}

No third BAE family, new BA2 operator, universal derivation DSL, general dependency graph or implementation schema is forced.

The active closure artifact is:

\begin{itemize}
\tightlist
\item
  \texttt{methodology/BA3\_PROVENANCE\_DERIVATION\_LIFECYCLE\_CHANGE\_CONTRACT\_R1.md}
\end{itemize}

The BA3-T1/T2/T3 artifacts remain immutable research derivation history.

\subsection{3. Base Analysis status after BA3}\label{base-analysis-status-after-ba3}

\begin{verbatim}
BA0   CLOSED
BA1   CLOSED
BA2   CLOSED
BA3   CLOSED
BA4   NOT STARTED / NEXT
BA5   NOT STARTED
BA6   NOT STARTED
\end{verbatim}

The Base Analysis milestone as a whole is not yet closed. BA4 must establish projection semantics, BA5 lexical/assistance boundaries and BA6 integrated regression/holdout closure.

\subsection{4. BA4 objective - projections}\label{ba4-objective---projections}

BA4 must define the smallest contract by which one accepted Base Analysis can support multiple derived human and methodology-oriented views without making any view a new source of project truth or allowing a consumer to redefine the shared core.

BA4 is not a system-modeling notation design and is not a formal STRIDE schema exercise.

Key questions include:

\begin{itemize}
\tightlist
\item
  what a projection may select, rename, group or aggregate from accepted BA semantics;
\item
  what projection-local identity/traceability is needed back to BA elements;
\item
  how a method-owned projection may add method-specific interpretation without feeding that taxonomy back into BA;
\item
  which information must remain directly recoverable from BA so a consumer does not reconstruct project meaning from raw prose;
\item
  how projection rebuild/revalidation follows BA change without becoming another lifecycle authority.
\end{itemize}

\subsection{5. BA4-T1 - projection boundary, traceability and semantic-preservation lower-bound trial}\label{ba4-t1---projection-boundary-traceability-and-semantic-preservation-lower-bound-trial}

\textbf{NEXT / NOT EXECUTED BY THIS PLAN REVISION.}

BA4-T1 must perform only the first projection-boundary trial.

It must pressure-test at least:

\begin{enumerate}
\def\labelenumi{\arabic{enumi}.}
\tightlist
\item
  one human-oriented projection and one bounded method-oriented projection derived from the same accepted BA materialization;
\item
  selection/renaming/grouping/aggregation that preserves BA meaning versus projection-owned interpretation that must remain downstream;
\item
  the minimum trace from a projection item/result back to one or more \texttt{BAReferent}/\texttt{BAProposition} identities;
\item
  whether projection items require any new first-class BAE identity family or can remain view/method-owned constructs;
\item
  whether a consumer can obtain the needed shared project meaning without re-reading governed prose or importing method-specific semantics into BA;
\item
  how a BA identity replacement/retirement affects a projection without creating a second project lifecycle model;
\item
  whether any counterexample forces the smallest BA1/BA2/BA3 reopen.
\end{enumerate}

\subsubsection{BA4-T1 guardrails}\label{ba4-t1-guardrails}

Do not:

\begin{itemize}
\tightlist
\item
  define a complete STRIDE or STRIDE-AI ontology/schema;
\item
  define AnalysisRecord or Common Finding;
\item
  design ThreatForge classes/tables/UI;
\item
  make a projection project authority;
\item
  make DFD elements first-class BA types merely because one method consumes them;
\item
  reopen Chapters 2-4 or BA0-BA3 without a material counterexample.
\end{itemize}

\subsubsection{BA4-T1 exit condition}\label{ba4-t1-exit-condition}

Produce a falsifiable projection responsibility/traceability lower-bound candidate and the next smallest pressure target. Do not close BA4 from one projection example.

\subsection{6. BA5 - lexical vocabulary and optional assistance}\label{ba5---lexical-vocabulary-and-optional-assistance}

Only after BA4 boundaries are stable, define display labels, authoring synonyms and optional assistance. Preserve:

\begin{verbatim}
source lexical wording
        !=
stable BA semantic key
        !=
display label / authoring synonym
\end{verbatim}

NLP/LLM support remains candidate-producing assistance subject to provenance/review; it does not become semantic authority.

\subsection{7. BA6 - complete Base Analysis regression and closure}\label{ba6---complete-base-analysis-regression-and-closure}

Regress the complete BA0-BA5 design against the closed corpora and at least one structurally different holdout. BA6 may reopen only the smallest earlier responsibility on a material counterexample.

Only BA6 may close the complete Base Analysis milestone for the current thesis scope.

\subsection{8. Analysis envelope remains downstream}\label{analysis-envelope-remains-downstream}

Only after Base Analysis closure:

\begin{itemize}
\tightlist
\item
  A1 - AnalysisRecord / execution envelope;
\item
  A2 - common reviewed Finding boundary;
\item
  A3 - change/provenance integration;
\item
  formal methodology overlays and STRIDE/STRIDE-AI evaluations.
\end{itemize}

\subsection{9. Next authorized microstep}\label{next-authorized-microstep}

Execute only:

\begin{quote}
\textbf{BA4-T1 - projection boundary, traceability and semantic-preservation lower-bound trial.}
\end{quote}

Do not start BA5, BA6 or downstream analysis schemas before BA4-T1 is completed and reviewed.

\end{document}
